\begin{definition}[Derivative at a Point]
Let $f: D \to \mathbb{R}$ be a real-valued function and $c \in \text{int}(D)$. We say $f$ is differentiable at $c$ with derivative $L = f'(c)$ if and only if:
\[
\forall \epsilon \in \mathbb{R} \left( \epsilon > 0 \implies \exists \delta \in \mathbb{R} \left( \delta > 0 \land \forall x \in D \left( 0 < |x - c| < \delta \implies \left| \frac{f(x) - f(c)}{x - c} - L \right| < \epsilon \right) \right) \right)
\]
\end{definition}



\begin{definition}[Topological Definition of Derivative at a Point]
Let $f: D \to \mathbb{R}$ and $c \in \text{int}(D)$. $f$ is differentiable at $c$ with derivative $L$ if for every neighborhood $V$ of $L$, there exists a neighborhood $U$ of $c$ such that for all $x \in (U \setminus \{c\}) \cap D$:
\[
\frac{f(x) - f(c)}{x - c} \in V
\]
\end{definition}




\begin{definition}[Sequential Definition of Derivative at a Point]
Let $A \subseteq \mathbb{R}$, let $f : A \to \mathbb{R}$, and let $c \in A$ be a limit point of $A$. We say that $f$ is differentiable at $c$ with derivative $L$ if
\[
\forall (x_n) \subseteq A \setminus \{c\}\;
\Big(
x_n \to c
\;\Rightarrow\;
\frac{f(x_n) - f(c)}{x_n - c} \to L
\Big).
\]

\end{definition}





\begin{theorem}[Equivalence of Definitions of the Derivative at a Point]
\label{thm:derivative-equivalence}
Let $A \subseteq \mathbb{R}$, let $f : A \to \mathbb{R}$, and let $c \in A$ be a limit point of $A$. Let $L \in \mathbb{R}$. The following are equivalent:

\begin{enumerate}[label=(\roman*)]

\item \textbf{(Epsilon--Delta Definition)}
\[
\forall \epsilon > 0\;\exists \delta > 0\;\forall x \in A\;
\Big(
0 < |x - c| < \delta
\;\Rightarrow\;
\left|
\frac{f(x) - f(c)}{x - c} - L
\right| < \epsilon
\Big).
\]

\item \textbf{(Neighborhood Definition)}
For every $\epsilon > 0$, there exists a neighborhood $V$ of $c$ such that
\[
x \in (A \cap V) \setminus \{c\}
\;\Rightarrow\;
\left|
\frac{f(x) - f(c)}{x - c} - L
\right| < \epsilon.
\]

\item \textbf{(Sequential Definition)}
\[
\forall (x_n) \subseteq A \setminus \{c\}\;
\Big(
x_n \to c
\;\Rightarrow\;
\frac{f(x_n) - f(c)}{x_n - c} \to L
\Big).
\]

\end{enumerate}
\end{theorem}

\begin{remark}[The Principle of Controlled Approximation]
Across all three definitions, the underlying logical structure is the \textbf{control of the difference quotient}.

\begin{itemize}
    \item \textbf{Metric Control ($\varepsilon$--$\delta$):}
    Precision of the difference quotient relative to $L$ is regulated by the radius $\delta$ of the punctured domain.

    \item \textbf{Topological Control (Neighborhoods):}
    Containment of the difference quotient within a target set $V$ is guaranteed by the existence of a sufficiently small neighborhood $U$ of $c$.

    \item \textbf{Sequential Control (Sequences):}
    Any probe sequence $(x_n)$ converging to $c$ forces the induced sequence of difference quotients to converge to $L$.
\end{itemize}

\medskip

In applied settings (e.g., numerical simulation), this expresses a stability principle: if the input resolution (such as a time-step $\Delta t$) is sufficiently controlled, then the resulting approximation error in the difference quotient remains uniformly bounded. This provides a formal mathematical constraint underlying numerical stability.
\end{remark}

\begin{proposition}[Derivative: $h$-Form Equivalence]
Let $I \subseteq \mathbb{R}$ be an open interval, let $f : I \to \mathbb{R}$, and let $c \in I$. The following are equivalent:

\begin{enumerate}
\item $f$ is differentiable at $c$.

\item The limit
\[
\lim_{h \to 0} \frac{f(c + h) - f(c)}{h}
\]
exists.
\end{enumerate}

In this case,
\[
f'(c) = \lim_{h \to 0} \frac{f(c + h) - f(c)}{h}.
\]
\end{proposition}

\begin{remark}[The two formulations are equivalent]
Setting $h = x - c$ converts one form to the other bijectively.
The condition $0 < |x - c| < \delta$ becomes $0 < |h| < \delta$;
$f$ being undefined at $h = 0$ is already excluded by $h \neq 0$.
\end{remark}

\begin{remark}[Standard notation]
$f'(c)$, $\;\dfrac{df}{dx}\big|_{x=c}$, $\;Df(c)$, and $\;\dot{f}(c)$
all denote the same quantity.
Leibniz notation $\dfrac{df}{dx}$ makes the chain rule visually transparent.
Prime notation $f'$ is most compact for real analysis.
\end{remark}

\begin{remark}[What the definition requires]
Three things are needed for $f'(c)$ to exist:
\begin{enumerate}
  \item $f$ is defined on an open interval containing $c$,
        so the difference quotient is defined for $x$ near $c$ on both sides.
  \item The left and right limits of the difference quotient as $x \to c$ agree.
  \item That common limit is finite.
\end{enumerate}
If the limit is $\pm\infty$, the function has a vertical tangent at $c$
and is not differentiable there in this sense.
\end{remark}

\begin{remark}[Canonical non-example: $|x|$ at $0$]
For $f(x) = |x|$:
\[
\frac{f(x) - f(0)}{x - 0} = \frac{|x|}{x} =
\begin{cases} +1 & x > 0, \\ -1 & x < 0. \end{cases}
\]
The one-sided limits disagree, so $f'(0)$ does not exist.
This is the canonical example that continuity does not imply differentiability.
\end{remark}


















\begin{theorem}[Differentiable Implies Continuous]
\label{thm:differentiable-implies-continuous}
Let $A \subseteq \mathbb{R}$, let $f : A \to \mathbb{R}$, and let $c \in A$ be a limit point of $A$. If there exists $L \in \mathbb{R}$ such that $f'(c)=L$, then
\[
\lim_{x \to c} f(x) = f(c).
\]
Hence $f$ is continuous at $c$.
\end{theorem}





\begin{remark}[Converse fails]
Continuity does not imply differentiability.
$f(x) = |x|$ is continuous everywhere but fails to be differentiable at $0$.
More dramatically, the Weierstrass function is continuous everywhere
and differentiable nowhere.
\end{remark}

\begin{remark}[Position in the regularity hierarchy]
\[
\text{differentiable at } c
\;\Rightarrow\;
\text{continuous at } c
\;\not\Rightarrow\;
\text{differentiable at } c.
\]
Every assumption of differentiability is strictly stronger than an
assumption of continuity alone.
\end{remark}







\begin{theorem}[Uniqueness of the Derivative]
\label{thm:uniqueness-of-the-derivative}
Let $A \subseteq \mathbb{R}$, let $f : A \to \mathbb{R}$, and let $c \in A$ be a limit point of $A$. If $f$ is differentiable at $c$, then the derivative of $f$ at $c$ is unique.

Equivalently, if $L,M \in \mathbb{R}$ satisfy
\[
\forall \epsilon > 0\;\exists \delta > 0\;\forall x \in A\;
\Big(
0 < |x-c| < \delta
\;\Rightarrow\;
\left|
\frac{f(x)-f(c)}{x-c} - L
\right| < \epsilon
\Big),
\]
and
\[
\forall \epsilon > 0\;\exists \delta > 0\;\forall x \in A\;
\Big(
0 < |x-c| < \delta
\;\Rightarrow\;
\left|
\frac{f(x)-f(c)}{x-c} - M
\right| < \epsilon
\Big),
\]
then
\[
L=M.
\]
\end{theorem}


\begin{proposition}[Derivative of the Identity]
Let $I \subseteq \mathbb{R}$ be an open interval, and let $\operatorname{id}_I : I \to \mathbb{R}$ be the identity function defined by
\[
\operatorname{id}_I(x) := x.
\]
Then $\operatorname{id}_I$ is differentiable at every $c \in I$, and
\[
\operatorname{id}_I'(c) = 1.
\]
\end{proposition}